% Part: proof-theory
% Chapter: cut-elimination
% Section: ce-largest

\documentclass[../../../include/open-logic-section]{subfiles}

\begin{document}

\olfileid[fr]{pt}{cut}{inv}

\olsection{Élimination des coupures de rang maximal}

La réduction repose sur l'inversion des règles logiques avec
une borne sur le nombre de coupures de chaque rang.












La preuve du \olref[seq][inv]{lem:G3c-invert} montre que, lorsque
\Log{G3c} !!{prove}s la conclusion d'une règle logique autre que
\RightR{\lexists} ou \LeftR{\lforall}, il !!{prove}s aussi
chacune de ses prémisses sans augmenter la !!{height} de la preuve.
Nous allons étendre ce résultat à $\Log{G3c}+\CutCS$.

Comme précédemment, le \emph{rang} d'une inférence \CutCS est
la profondeur de sa !!{formula} de coupure.

\begin{lem}\ollabel{lem:inv-G3c-cut}
Si $\Log{G3c}+\CutCS$ !!{prove}s la conclusion d'une règle
logique~$R$, autre que \RightR{\lexists} ou \LeftR{\lforall},
au moyen d'une !!{proof}~$\pi$, il !!{prove}s aussi chacune de
ses prémisses au moyen d'une !!{proof}~$\pi'$ ou, pour une règle
à deux prémisses, de !!{proof}s~$\pi'$ et $\pi''$.
En outre,
(a)~$\pheight{\pi'}\le\pheight{\pi}$ et, le cas échéant,
$\pheight{\pi''}\le\pheight{\pi}$;
(b)~pour chaque rang, le nombre d'inférences \CutCS de ce rang
dans~$\pi'$ et, séparément, dans~$\pi''$, n'excède pas leur nombre
dans~$\pi$.
Pour les règles universelle à droite et existentielle à gauche,
la constante de la prémisse inversée doit être fraîche dans
le séquent conclusion, comme dans le lemme d'inversion quantifiée.
\end{lem}

\begin{proof}
Examinons les preuves des
\cref{pt:seq:inv:lem:G3c-invert,pt:seq:inv:lem:invert-quant}.
La récurrence porte sur~$\pheight{\pi}$.
Au cas initial, on remplace un axiome par un autre axiome;
les conditions (a) et~(b) sont donc immédiates.
Si la dernière règle de~$\pi$ est précisément la règle~$R$
que l'on inverse, ses sous-!!{proof}s immédiates donnent
les !!{proof}s~$\pi_i$ recherchées. Leur hauteur est plus
petite, et leurs coupures forment des sous-ensembles des
occurrences de coupure de la preuve entière. La condition
(b) est une borne supérieure, et non une égalité.
Dans les autres cas logiques, on applique l'hypothèse de
récurrence aux sous-!!{proof}s immédiates, puis on réapplique
la règle qui terminait la preuve. Les bornes de hauteur
et de nombre de coupures de chaque rang se transmettent ainsi
à la preuve reconstruite.
Dans le cas quantifié, les renommages de constantes nécessaires
ne changent pas la profondeur des formules de coupure.

Il reste le cas où la dernière inférence est \CutCS.
Traitons à nouveau un exemple : soit $\pi$ une !!{proof}
de la conclusion $!B\land!C,\Gamma\Sequent\Delta$
de \LeftR{\land}, qui se termine par
\[
\AxiomC{}
\RightLabel{$\pi_1$}
\Deduce$!B \land !C, \Gamma \fCenter \Delta, !A$
\AxiomC{}
\RightLabel{$\pi_2$}
\Deduce$!A, !B \land !C, \Gamma \fCenter \Delta$
\RightLabel{\CutCS}
\BinaryInf$!B \land !C, \Gamma \fCenter \Delta$
\DisplayProof
\]
L'hypothèse de récurrence s'applique à $\pi_1$ et~$\pi_2$,
qui démontrent aussi des instances de la conclusion de
$\LeftR{\land}$. Elle fournit des !!{proof}s
$\pi_1'$ et~$\pi_2'$ respectivement de
$!B,!C,\Gamma\fCenter\Delta,!A$ et
$!A,!B,!C,\Gamma\fCenter\Delta$.
On les réunit par \CutCS pour obtenir~$\pi'$ :
\[
\AxiomC{}
\RightLabel{$\pi_1'$}
\Deduce$!B, !C, \Gamma \fCenter \Delta, !A$
\AxiomC{}
\RightLabel{$\pi_2'$}
\Deduce$!A, !B, !C, \Gamma \fCenter \Delta$
\RightLabel{\CutCS}
\BinaryInf$!B, !C, \Gamma \fCenter \Delta$
\DisplayProof
\]
Les conditions (a) et~(b) sont satisfaites : la hauteur ne croît
pas, et l'on ajoute aux coupures des deux prémisses la même
coupure finale que dans la preuve initiale.
\end{proof}

% BEGIN EDITORIAL AUXILIARIES
\paragraph{Complément éditorial : règles auxiliaires en présence de coupures.}
Les arguments précédents d'affaiblissement, de substitution dans une
preuve régulière et de contraction s'étendent à
$\Log{G3c}+\CutCS$, sans introduire de coupure d'un rang supérieur
aux rangs déjà présents. Pour l'affaiblissement et la substitution,
on ajoute à la récurrence sur la hauteur le cas d'une dernière
règle \CutCS : on transforme ses deux prémisses et on réapplique
la même coupure. La substitution de termes ne change pas la
profondeur des formules.

Pour la contraction, on suit également la récurrence sur la hauteur
de la preuve. Dans les cas logiques où la formule contractée est
principale, on utilise le lemme d'inversion ci-dessus sur les
sous-preuves immédiates, puis l'hypothèse de récurrence, puis on
réapplique la règle logique. L'inversion n'augmente pas la hauteur
ni les rangs des coupures, de sorte que la récurrence porte toujours
sur des preuves de hauteur strictement plus petite. Si la dernière
règle est \CutCS, les deux occurrences à contracter appartiennent
à son contexte commun : on les contracte dans chacune des deux
prémisses par récurrence, puis on réapplique la même coupure.
On distingue ces occurrences de la formule coupée, même si les
formules sont identiques. Cet argument n'utilise pas le théorème
d'élimination des coupures.

Avant les transformations, on renomme les variables propres pour
rendre les preuves régulières et les tenir à l'écart des termes et
des contextes qui seront ajoutés. Ces renommages ne changent ni la
hauteur ni le rang des coupures.
% END EDITORIAL AUXILIARIES

Ce raisonnement ne fonctionnerait pas directement avec \Cut
à la place de \CutCS : le séquent final contiendrait deux copies
de $!B,!C,\Gamma$ à gauche et deux copies de $\Delta$ à droite.
Il faudrait utiliser l'admissibilité de la contraction; or la
preuve de la \olref[seq][inv]{prop:G3c-cont-adm} fait appel au lemme
d'inversion. Pour établir l'inversion de cette façon, l'argument
serait donc circulaire.

Le \cref{pt:cut:inv:lem:inv-G3c-cut} permet de donner une autre preuve
de l'élimination des coupures. Nous ne supprimons plus successivement
les inférences \CutCS les plus hautes de l'arbre, mais celles
de rang maximal. Partant d'une !!{proof}~$\pi$ dans
$\Log{G3c}+\CutCS$, nous supprimons une coupure, où qu'elle
se trouve, en la remplaçant éventuellement par des coupures
de rang inférieur. Nous répétons cette opération jusqu'à
ce qu'il ne reste plus de coupure. La seule condition locale
est qu'aucune coupure de rang égal ou supérieur ne se trouve
au-dessus de celle que nous traitons.

\begin{lem}\ollabel{lem:max-cut-red-G3c}
Si $\pi$ est une !!{proof} dans $\Log{G3c}+\CutCS$
qui se termine par une inférence \CutCS de rang~$n$ et dont
toutes les autres inférences sont des règles de~$\Log{G3c}$
ou des coupures \CutCS de rang $<n$, il existe une preuve~$\pi'$
du même séquent final dans~$\Log{G3c}+\CutCS$ dont toutes
les inférences \CutCS ont un rang~$<n$.
\end{lem}

\begin{proof}
La !!{proof}~$\pi$ se termine par
\[
\AxiomC{}
\RightLabel{$\pi_1$}
\Deduce$\Gamma \fCenter \Delta, !A$
\AxiomC{}
\RightLabel{$\pi_2$}
\Deduce$!A, \Gamma \fCenter \Delta$
\RightLabel{\CutCS}
\BinaryInf$\Gamma \fCenter \Delta$
\DisplayProof
\]
Nous distinguons plusieurs cas suivant la forme de~$!A$.

% BEGIN EDITORIAL FALSE
\paragraph{Complément éditorial : constante de fausseté.}
Si $!A\ident\lfalse$, la coupure a le rang zéro, et ses deux
prémisses ne contiennent aucune coupure. On obtient directement
une preuve de $\Gamma\Sequent\Delta$ en supprimant, dans toute
la première sous-preuve, la copie suivie de la constante de
fausseté à droite. Cette formule ne peut être principale dans une
règle logique ni constituer l'atome commun d'un axiome d'identité.
Les axiomes de fausseté à gauche restent des axiomes. Le résultat
est donc une preuve sans coupure. Ce cas est distinct du cas
atomique suivant.
% END EDITORIAL FALSE

Supposons que $!A$ soit atomique.
Toutes les inférences \CutCS de~$\pi_1$ et $\pi_2$ doivent
avoir un rang $<\depth{!A}=0$; ces deux preuves ne contiennent
donc aucune coupure.
Montrons par récurrence sur $\pheight{\pi_2}$ que
$\Gamma\Sequent\Delta$ possède une preuve sans coupure.
Si $\pheight{\pi_2}=0$, le séquent
$!A,\Gamma\Sequent\Delta$ est un axiome.
S'il s'agit d'un axiome de fausseté, celle-ci appartient déjà
au contexte gauche, puisque la formule coupée est atomique
et distincte de la constante de fausseté; le séquent recherché
est encore un axiome.
Pour un axiome d'identité, si $!A$ n'est pas principal,
un atome $!C$ appartient à la fois à~$\Gamma$ et à $\Delta$,
et $\Gamma\Sequent\Delta$ est lui-même un axiome.
Si $!A$ est principal, écrivons $\Delta=\Delta',!A$.
La preuve $\pi_1$ se termine alors par
$\Gamma\Sequent\Delta',!A,!A$.
L'admissibilité de la contraction,
\olref[seq][inv]{prop:G3c-cont-adm}, donne une preuve sans
coupure de $\Gamma\Sequent\Delta',!A$.

Si $\pheight{\pi_2}>0$, la preuve se termine par une règle~$R$.
Une prémisse de cette règle est de la forme
$!A,\Gamma',\Pi\Sequent\Delta',\Lambda$,
où $\Pi$ et $\Lambda$ sont les !!{formula}s actives de~$R$.
L'affaiblissement \olref[seq][adm]{prop:weak-G3c-adm},
appliqué à~$\pi_1$ et à la sous-preuve de~$\pi_2$ qui conduit
à cette prémisse, donne des !!{proof}s de
$\Gamma,\Gamma',\Pi\Sequent\Delta,\Delta',\Lambda,!A$ et
$!A,\Gamma,\Gamma',\Pi\Sequent\Delta,\Delta',\Lambda$.
La seconde a une hauteur strictement inférieure à celle
de la preuve~$\pi_2$ initiale, car l'affaiblissement conserve
la hauteur. L'hypothèse de récurrence fournit donc, pour
chaque prémisse de~$R$, une !!{proof} de
$\Gamma,\Gamma',\Pi\Sequent\Delta,\Delta',\Lambda$.
En réappliquant $R$, on obtient
$\Gamma,\Gamma\Sequent\Delta,\Delta$.
Enfin, \olref[seq][inv]{prop:G3c-cont-adm} donne la preuve
sans coupure recherchée de $\Gamma\Sequent\Delta$.

Si $!A$ n'est ni atomique ni la constante de fausseté,
nous distinguons les cas suivant son !!{main operator}.
Dans chaque cas où les règles correspondantes sont inversibles,
\olref{lem:inv-G3c-cut} fournit les !!{proof}s de leurs prémisses
avec $!A$ comme formule principale.
Nous procédons ensuite comme dans le cas~E de la preuve
du \olref[top]{lem:cut-adm-G3c}.
\begin{enumerate}
  \item $!A\ident\lnot!B$.
  Les !!{proof}s $\pi_1$ et $\pi_2$ se terminent par
  $\Gamma\Sequent\Delta,\lnot!B$ et
  $\lnot!B,\Gamma\Sequent\Delta$.
  D'après le \olref{lem:inv-G3c-cut}, il existe des !!{proof}s
  $\pi_1'$ et $\pi_2'$ de $!B,\Gamma\Sequent\Delta$
  et de $\Gamma\Sequent\Delta,!B$.
  La !!{proof}~$\pi'$ est
  \[
  \AxiomC{}
  \RightLabel{$\pi_2'$}
  \Deduce$\Gamma \fCenter \Delta, !B$
  \AxiomC{}
  \RightLabel{$\pi_1'$}
  \Deduce$!B, \Gamma \fCenter \Delta$
  \RightLabel{\CutCS}
  \BinaryInf$\Gamma \fCenter \Delta$
  \DisplayProof
  \]
  \item $!A\ident(!B\lor!C)$.
  Les !!{proof}s $\pi_1$ et $\pi_2$ se terminent par
  $\Gamma\Sequent\Delta,!B\lor!C$ et
  $!B\lor!C,\Gamma\Sequent\Delta$.
  L'inversion de \RightR{\lor} dans~$\pi_1$,
  par le \olref{lem:inv-G3c-cut}, fournit une !!{proof} $\pi_1'$
  de $\Gamma\Sequent\Delta,!B,!C$.
  L'inversion de \LeftR{\lor} dans~$\pi_2$,
  par le \olref{lem:inv-G3c-cut}, fournit une !!{proof} $\pi_2'$
  de $!B,\Gamma\Sequent\Delta$ et une !!{proof} $\pi_2''$
  de $!C,\Gamma\Sequent\Delta$.
  L'affaiblissement \olref[seq][adm]{prop:weak-G3c-adm}
  appliqué à $\pi_2''$ donne aussi une !!{proof}~$\pi_2'''$
  de $!C,\Gamma\Sequent\Delta,!B$.
  La !!{proof}~$\pi'$ est
  \[
  \AxiomC{}
  \RightLabel{$\pi_1'$}
  \Deduce$\Gamma \fCenter \Delta, !B, !C$
  \AxiomC{}
  \RightLabel{$\pi_2'''$}
  \Deduce$!C, \Gamma \fCenter \Delta, !B$
  \RightLabel{\CutCS}
  \BinaryInf$\Gamma \fCenter \Delta, !B$
  \AxiomC{}
  \RightLabel{$\pi_2'$}
  \Deduce$!B, \Gamma \fCenter \Delta$
  \RightLabel{\CutCS}
  \BinaryInf$\Gamma \fCenter \Delta$
  \DisplayProof
  \]
  \item $!A\ident(!B\land!C)$. Exercice.
  \item $!A\ident(!B\lif!C)$.
  Les !!{proof}s $\pi_1$ et $\pi_2$ se terminent par
  $\Gamma\Sequent\Delta,!B\lif!C$ et
  $!B\lif!C,\Gamma\Sequent\Delta$.
  L'inversion de \RightR{\lif} dans~$\pi_1$,
  par le \olref{lem:inv-G3c-cut}, fournit une !!{proof} $\pi_1'$
  de $!B,\Gamma\Sequent\Delta,!C$.
  L'inversion de \LeftR{\lif} dans~$\pi_2$,
  par le \olref{lem:inv-G3c-cut}, fournit une !!{proof} $\pi_2'$
  de $\Gamma\Sequent\Delta,!B$ et une !!{proof} $\pi_2''$
  de $!C,\Gamma\Sequent\Delta$.
  L'affaiblissement \olref[seq][adm]{prop:weak-G3c-adm},
  appliqué à $\pi_2''$, donne une !!{proof}~$\pi_2'''$
  de $!C,!B,\Gamma\Sequent\Delta$.
  La !!{proof}~$\pi'$ est
  \[
  \AxiomC{}
  \RightLabel{$\pi_2'$}
  \Deduce$\Gamma \fCenter \Delta, !B$
  \AxiomC{}
  \RightLabel{$\pi_1'$}
  \Deduce$!B, \Gamma \fCenter \Delta, !C$
  \AxiomC{}
  \RightLabel{$\pi_2'''$}
  \Deduce$!C, !B, \Gamma \fCenter \Delta$
  \RightLabel{\CutCS}
  \BinaryInf$!B, \Gamma \fCenter \Delta$
  \RightLabel{\CutCS}
  \BinaryInf$\Gamma \fCenter \Delta$
  \DisplayProof
  \]
  \item $!A\ident\lforall[x][!B(x)]$.
  Les !!{proof}s $\pi_1$ et $\pi_2$ se terminent par
  $\Gamma\Sequent\Delta,\lforall[x][!B(x)]$ et
  $\lforall[x][!B(x)],\Gamma\Sequent\Delta$.
  Par le \olref{lem:inv-G3c-cut}, on obtient une !!{proof}~$\pi_1'$
  de $\Gamma\Sequent\Delta,!B(c)$.
  On choisit la constante de cette inversion fraîche dans
  les deux preuves initiales. On rend la preuve obtenue
  régulière en renommant ses variables propres afin qu'elles
  n'apparaissent dans aucun terme ni contexte de l'autre preuve.
  Les substitutions et affaiblissements ci-dessous respectent
  alors toutes les conditions de variables propres.

  Examinons maintenant la !!{proof}~$\pi_2$, dont le séquent
  final est $\lforall[x][!B(x)],\Gamma\Sequent\Delta$.
  En remontant l'arbre à partir de ce séquent, marquons les
  occurrences de $\lforall[x][!B(x)]$ à gauche qui conduisent
  à l'occurrence coupée. Plus précisément :
  (a)~marquons cette occurrence dans le séquent final de~$\pi_2$;
  (b)~lorsqu'une occurrence marquée se trouve dans le contexte
  de la conclusion d'une règle, marquons l'occurrence
  correspondante dans chacune de ses prémisses;
  (c)~si l'occurrence marquée est principale dans la conclusion
  d'une règle \LeftR{\lforall}, marquons les occurrences
  actives correspondantes de $\lforall[x][!B(x)]$ et de $!B(t)$
  dans sa prémisse.

  Considérons toutes les occurrences ainsi marquées de
  $\lforall[x][!B(x)]$. Aucune n'est active pour une règle autre
  que \LeftR{\lforall}. Supprimons ces occurrences dans~$\pi_2$.
  Si le résultat reste une !!{proof} correcte, ces occurrences
  n'étaient jamais actives et provenaient directement des
  axiomes. La preuve se termine alors par
  $\Gamma\Sequent\Delta$; elle remplace $\pi$ et supprime
  la coupure de rang maximal.

  Sinon, certaines !!{formula}s supprimées étaient actives
  dans des inférences \LeftR{\lforall}.
  Ces inférences sont devenues des passages incorrects
  de la forme
  \[
  \AxiomC{}
  \RightLabel{$\pi_t$}
  \Deduce$!B(t), \Pi \fCenter \Lambda$
  \RightLabel{\LeftR{\lforall}}
  \UnaryInf$\Pi \fCenter \Lambda$
  \DisplayProof
  \]
  Remplaçons chaque passage incorrect le plus haut par
  \[
  \AxiomC{}
  \RightLabel{$\Subst{\pi_1'}{t}{c}[\Pi \Sequent \Lambda]$}
  \Deduce$\Gamma, \Pi \fCenter \Delta, \Lambda, !B(t)$
  \AxiomC{}
  \RightLabel{$\pi_t[\Gamma \Sequent \Delta]$}
  \Deduce$!B(t), \Gamma, \Pi \fCenter \Delta, \Lambda$
  \RightLabel{\CutCS}
  \BinaryInf$\Gamma, \Pi \fCenter \Delta, \Lambda$
  \DisplayProof
  \]
  et ajoutons $\Gamma$ à gauche et $\Delta$ à droite de
  chaque séquent situé au-dessous.
  Aux règles à deux prémisses, on affaiblit aussi l'autre
  sous-preuve pour que les deux prémisses conservent des
  contextes communs. On procède de même lorsqu'on traverse
  une coupure déjà présente.
  Le passage le plus haut a une sous-preuve correcte
  au-dessus de sa prémisse; les nouvelles prémisses
  proviennent donc bien de preuves, par substitution
  régulière et affaiblissement.
  Répétons l'opération jusqu'à ce que toutes les anciennes
  inférences \LeftR{\lforall} dont la prémisse contient
  une occurrence marquée de $!B(t)$ soient remplacées
  par une coupure \CutCS sur~$!B(t)$.
  La !!{proof} obtenue est correcte, et son séquent final
  est de la forme
  $\Gamma,\dots,\Gamma\Sequent\Delta,\dots,\Delta$.
  Par la contraction en présence de coupures justifiée
  plus haut, on obtient une !!{proof} de
  $\Gamma\Sequent\Delta$ qui remplace $\pi$.
  La coupure sur $\lforall[x][!B(x)]$ a été remplacée
  par des coupures, éventuellement nombreuses, sur $!B(t)$,
  toutes de rang inférieur. Les coupures déjà présentes
  dans $\pi_1$ ou $\pi_2$ peuvent être recopiées au cours
  du procédé, mais leur rang reste lui aussi inférieur.
% BEGIN EDITORIAL EXISTS
  \item \emph{Complément éditorial : cas existentiel, omis dans la source.}
  Si $!A\ident\lexists[x][!B(x)]$, on procède de façon duale.
  L'inversion de la règle existentielle à gauche appliquée à
  $\pi_2$ fournit une preuve $\pi_2'$ de
  $!B(c),\Gamma\Sequent\Delta$, avec une constante $c$ fraîche.
  On suit, dans $\pi_1$, les occurrences de la formule existentielle
  à droite qui conduisent à l'occurrence coupée.
  À chaque règle existentielle à droite portant sur une occurrence
  suivie, on suit aussi la formule active $!B(t)$ dans la prémisse.
  On supprime les occurrences suivies de la formule existentielle.

  Si le résultat est déjà une preuve correcte, il fournit le
  séquent recherché. Sinon, chaque ancienne application
  existentielle à droite dont la formule active est suivie a
  désormais pour prémisse $\Pi\Sequent\Lambda,!B(t)$ et pour
  conclusion $\Pi\Sequent\Lambda$. En commençant par la plus
  haute de ces applications, on la remplace par la coupure
  \[
  \AxiomC{}
  \RightLabel{$\pi_t[\Gamma\Sequent\Delta]$}
  \Deduce$\Gamma,\Pi\fCenter\Delta,\Lambda,!B(t)$
  \AxiomC{}
  \RightLabel{$\Subst{\pi_2'}{t}{c}[\Pi\Sequent\Lambda]$}
  \Deduce$!B(t),\Gamma,\Pi\fCenter\Delta,\Lambda$
  \RightLabel{\CutCS}
  \BinaryInf$\Gamma,\Pi\fCenter\Delta,\Lambda$
  \DisplayProof
  \]
  Ici, la première prémisse provient de la sous-preuve déjà
  corrigée par affaiblissement, et la seconde de la substitution
  régulière suivie d'affaiblissement. La constante fraîche et
  les renommages sont choisis comme dans le cas universel.
  On ajoute les contextes communs à tous les séquents situés
  au-dessous de cette coupure, en affaiblissant également
  les autres branches aux règles à deux prémisses pour
  conserver les contextes communs, puis on répète le procédé.
  La contraction finale supprime les répétitions des contextes.
  Toutes les nouvelles coupures portent sur $!B(t)$ et ont donc
  un rang strictement inférieur à celui de la formule existentielle.
  Les coupures déjà présentes conservent un rang inférieur.
% END EDITORIAL EXISTS
\end{enumerate}
\end{proof}

\begin{prob}
Vérifier le cas $!A\ident(!B\land!C)$ dans la preuve du
\olref{lem:max-cut-red-G3c}.
Autrement dit, montrer que, s'il existe une !!{proof}~$\pi$
qui se termine par
\[
\AxiomC{}
\RightLabel{$\pi_1$}
\Deduce$\Gamma \fCenter \Delta, !B \land !C$
\AxiomC{}
\RightLabel{$\pi_2$}
\Deduce$!B \land !C, \Gamma \fCenter \Delta$
\RightLabel{\CutCS}
\BinaryInf$\Gamma \fCenter \Delta$
\DisplayProof
\]
où $\pi_1$ et $\pi_2$ ne contiennent que des coupures
de rang $<\depth{!B\land!C}$, il existe une !!{proof}~$\pi'$
de $\Gamma\Sequent\Delta$ ne contenant que des coupures
de rang $<\depth{!B\land!C}$.
\end{prob}

% \begin{prob}
% Vérifier le cas $!A \ident \lexists[x][!B(x)]$ dans la preuve de
% \olref{lem:inv-G3c-cut}.
% \end{prob}

Le \olref{lem:max-cut-red-G3c} montre qu'une inférence \CutCS
de rang maximal peut être remplacée par des inférences \CutCS
de rang inférieur, pourvu qu'aucune coupure de même rang
ne se trouve au-dessus d'elle. En appliquant successivement
ce lemme aux sous-!!{proof}s qui se terminent par une coupure
de rang maximal la plus haute, on peut éliminer toutes
les inférences \CutCS d'une !!{proof}.

\begin{thm}\ollabel{thm:cut-elim-G3c-largest}
Toute !!{proof} $\pi$ dans $\Log{G3c}+\CutCS$ peut être
transformée en une preuve dans~\Log{G3c} du même séquent final.
\end{thm}

\begin{proof}
Si la preuve ne contient aucune coupure, il n'y a rien à faire.
Sinon, montrons d'abord que toutes ses coupures de rang maximal
peuvent être supprimées.
Soit $d$ ce rang maximal et $n$ le nombre de coupures de rang~$d$.
Raisonnons par récurrence sur~$n$.
Si $n=0$, le résultat est immédiat.
Si $n>0$, choisissons une coupure de rang~$d$ la plus haute,
et soit $\pi_1$ la sous-!!{proof} qui se termine par elle.
D'après le \olref{lem:max-cut-red-G3c}, il existe une
!!{proof}~$\pi_1'$ du même séquent final dont toutes les
coupures ont un rang $<d$. Remplaçons $\pi_1$ par $\pi_1'$
dans~$\pi$. La nouvelle !!{proof} contient $n-1$ coupures
de rang~$d$; l'hypothèse de récurrence s'applique.

Le résultat global suit par récurrence sur~$d$.
Si $d=0$, toutes les coupures ont le rang zéro, et l'étape
précédente les supprime toutes.
Si $d>0$, cette étape fournit une preuve où toutes les
coupures de rang $d$ ont été remplacées par des coupures
de rang~$<d$. Si aucune coupure ne reste, la transformation
est terminée. Sinon, le rang maximal restant est strictement
inférieur à celui des coupures de~$\pi$, et l'hypothèse
de récurrence s'applique.
\end{proof}

\end{document}
